% ==================== 1 引言 ====================
\section{引　言}

X 射线衍射（X-ray Diffraction，XRD）是研究晶体结构最主要的实验手段之一\cite{liang1991xray}。
当 X 射线照射到晶体上时，晶格中的原子成为次级波源，各方向散射波发生干涉，在满足布拉格
方程 $2d_{hkl}\sin\theta=n\lambda$ 的方向上产生衍射。衍射峰的位置（$2\theta$）由晶面间距
$d_{hkl}$ 决定，反映晶胞的形状与大小；衍射峰的强度由结构因子 $F_{hkl}$ 决定，反映晶胞内
原子的种类、坐标与占有率。因此衍射图谱常被比作晶体结构的``指纹''：它既能用于物相鉴定、
晶格参数测定、晶粒尺寸与微观应力分析，也可结合结构解析与 Rietveld 精修，从衍射数据反推
出三维晶体结构。

要读懂这些信息，前提是掌握晶体的描述方法、能够正确地建立并观察结构模型。晶体结构 =
点阵 + 结构基元：点阵描述平移周期性，空间群描述对称性，晶胞内原子坐标描述基元内容。
建立起正确的结构模型，才有可能理解各衍射峰的指标化、系统消光规律与强度变化，也才能为
后续第一性原理计算提供初始结构\cite{huang1988solid}。

本文内容是 X 射线晶体学课程两次建模作业的整理与补充。两次作业使用的 Diamond 软件
（Crystal Impact 公司产品\cite{diamond2023manual}）只需输入空间群、晶格常数与原子坐标
即可生成三维结构模型，也能模拟粉末衍射谱、显示晶面与配位多面体。第一次作业用它依次
建立了面心立方（FCC）、体心立方（BCC）、六方密堆积（HCP）、氯化钠（NaCl）、金红石、
锐钛矿、萤石与金刚石八种晶体结构，并逐项讲解其结构特征；第二次作业用同一软件演示
反演中心、旋转轴、旋转反演轴、镜面、螺旋轴与滑移面六类对称元素（见 2.4 节），对比
单胞与超晶胞中等效原子的分布，以弄清对称元素的作用方式及其与空间群符号、等效位置
的关系。两次作业合并成文，并补充了软件操作流程与课件中与建模直接相关的晶体学基础
（第 2.1、2.2 节）。
